\documentclass[11pt]{article}

% Page layout
\usepackage[a4paper,margin=1in]{geometry}
\usepackage{hyperref}

% Fonts and formatting
\usepackage{times}
\usepackage{setspace}
\usepackage{titlesec}

% Math and figures
\usepackage{amsmath,amssymb}
\usepackage{graphicx}

% Citations and links
\usepackage{cite}
\usepackage{hyperref}

% Section style
\titleformat{\section}{\large\bfseries}{\thesection}{1em}{}
\titleformat{\subsection}{\normalsize\bfseries}{\thesubsection}{1em}{}

% Line spacing
\setstretch{1.1}

% Title
\title{\textbf{Graphormer - CS768 Course Project}}

\author{
Harshil Amitkumar Solanki (23B1016),  Sumedh S S (23B1079) \\
}

\date{}

\begin{document}

\maketitle

%------------------------------------------------

\begin{abstract}
Graph Neural Networks (GNNs) have achieved strong performance in graph
representation learning tasks such as node classification, link prediction,
and graph-level prediction. However, most existing GNN architectures assume
fixed graph structures and homogeneous node and edge types, which limits
their effectiveness when dealing with misspecified or heterogeneous graphs.

Graphormer\cite{graphormer}, on the other hand, adapts the standard Transformer architecture
for graph representation learning by incorporating structural encodings
that capture key properties of graph topology. These encodings enable the
model to effectively utilize global attention while preserving graph
structure information. Furthermore, the expressive power of Graphormer has
been shown to subsume several popular GNN variants as special cases.

We plan to study this approach, benchamrk it with other SOTA methods and try to perform an ablation study while trying to extend the architecture.Our code is available at \url{https://github.com/HARSHIL2836R/Graphormer-Project.git}.
\end{abstract}
%------------------------------------------------

\section{Introduction}

Graphs are a fundamental data structure for representing relational information
in many real-world domains, including social networks, citation networks,
biological systems, and recommender systems. Graph Neural Networks (GNNs) have emerged as a powerful
approach for this task and have achieved strong performance in applications
such as node classification, link prediction, and graph classification.

Despite their success, most existing GNN models rely on a fixed graph structure
and typically assume homogeneous node and edge types. This assumption limits
their effectiveness when the input graph is noisy, incomplete, or heterogeneous.
Meanwhile, Transformer architectures have demonstrated  success in
modeling sequential and structured data across domains such as natural language
processing, computer vision, and speech processing.
Graphormer adapts the
standard Transformer architecture for graph representation learning by
introducing structural encoding techniques that incorporate node centrality,
spatial relations, and edge features into the attention mechanism.

%------------------------------------------------

% \section{Problem Formulation}

% Clearly define the problem your project aims to solve.

% Describe:

% \begin{itemize}
% \item Input data
% \item Desired outputs
% \item Constraints and assumptions
% \item Objective of the system
% \end{itemize}

% If applicable, include a mathematical formulation:

% \begin{equation}
% y = f(x; \theta)
% \end{equation}

% where $x$ represents the input and $\theta$ represents the model
% parameters.

\section{Problem Formulation}

Let a graph be defined as $G = (V, E)$ where $V$ is the set of nodes and $E \subseteq V \times V$ is the set of edges. Each node $v \in V$ has a feature vector $x_v \in \mathbb{R}^d$. The node features of the entire graph are represented as $X \in \mathbb{R}^{|V| \times d}$. The graph structure can be represented using an adjacency matrix $A \in {0,1}^{|V| \times |V|}$.

The goal of graph representation learning is to learn a parameterized function $f_\theta(G, X)$ that produces useful node representations for downstream tasks. For node-level prediction, the model computes embeddings

$$
h_v = f_\theta(v, G, X)
$$

which are used to predict node labels

$$
y_v = g(h_v)
$$

where $g(\cdot)$ is typically a classifier.

Traditional Graph Neural Networks (GNNs) compute node embeddings through neighborhood aggregation. At layer $k$, the representation of node $v$ is updated as

$$
h_v^{(k)} = \sigma \left(
W^{(k)} \cdot
\text{AGG}\left({h_u^{(k-1)} : u \in \mathcal{N}(v)}\right)
\right)
$$

where $\mathcal{N}(v)$ denotes the neighbors of node $v$, $W^{(k)}$ are learnable parameters, and $\sigma$ is a nonlinear activation function.

Transformer-based models instead use global self-attention, allowing each node to attend to every other node. Given node embeddings $H$, attention is computed as

$$
\text{Attention}(Q,K,V) =
\text{softmax}\left(\frac{QK^T}{\sqrt{d_k}}\right)V
$$

where $Q = HW_Q$, $K = HW_K$, and $V = HW_V$ are projections with learnable parameters.

However, standard transformers do not explicitly encode graph structure. The central problem studied in this project is whether transformers inherently perform poorly for graph representation learning, or whether their performance can match or exceed GNNs when appropriate structural encodings are incorporated.


%------------------------------------------------

\section{Literature Review}

\subsection{Graph Transformer Networks \cite{gtn}}

Traditional Graph Neural Networks (GNNs) assume that the graph structure is fixed. 
During training, a GNN only aggregates information from the neighbors already 
connected in the given adjacency matrix.
Graph Transformer Networks (GTNs) address this limitation by learning new graph 
connections automatically. Instead of relying on a single adjacency matrix, GTNs 
start with multiple adjacency matrices corresponding to different edge types 
(e.g., author–paper, paper–conference). The model then learns weights that softly 
select and combine these edge types. By multiplying selected adjacency matrices, 
the model effectively creates multi-hop relations between nodes. Such paths are called \emph{meta-paths}. 

Through this process, GTNs can discover useful indirect relationships and 
construct new graph structures that better capture dependencies between nodes. 
These learned graphs are then used by standard graph convolution layers to learn 
node representations.

\subsection{Path-Augmented Graph Transformer Networks \cite{pagtn}}

Path-Augmented Graph Transformer Networks (PAGTN) extend the Transformer architecture to graph-structured data by explicitly incorporating path-based structural information between node pairs. GCNs require stacking many layers for capturing long range dependencies, which may lead to information loss or oversmoothing. PAGTN addresses this limitation by enabling global attention across all nodes in the graph

To capture structural relationships between nodes, PAGTN introduces \emph{path features}. For every pair of nodes, the model computes the shortest path between them and extracts useful information from this path. These features may include the types of edges along the path, the distance between the nodes, and whether the nodes belong to the same ring structure (which is particularly important in molecular graphs).

During the attention process, the model considers both the node features and the path features when determining how strongly two nodes should influence each other. The updated node representation is then computed by combining information from all other nodes in the graph, weighted by the attention scores. In this way, PAGTN allows each node to gather information from the entire graph while still using the path features to understand the structural context of these interactions.

\subsection{GRPE: Relative Positional Encoding for Graph Transformers \cite{grpe}}

GRPE extends the idea of structural encodings in graph transformers by introducing relative positional encodings that capture structural relationships between node pairs. Instead of relying solely on shortest-path distance, the method learns embeddings that represent relative positions between nodes in the graph.

The work emphasizes that spatial relationships between nodes are critical for effective graph transformers. By incorporating relative positional information into the attention mechanism, GRPE improves the model's ability to capture structural dependencies across the graph.

This work builds upon the core idea introduced in Graphormer that transformers require explicit structural encodings in order to perform well on graph representation tasks.


\subsection{Structure-Aware Transformer (SAN) \cite{san}}

Structure-Aware Transformer (SAN) investigates the role of structural positional encodings in transformer-based graph models. The paper proposes incorporating spectral information derived from the graph Laplacian to represent node positions within the graph structure.

These spectral encodings provide global structural information that helps the attention mechanism distinguish different structural roles of nodes. The results demonstrate that structural positional information significantly improves the expressiveness of graph transformers.

SAN further supports the hypothesis presented in Graphormer that transformers can be effective for graph learning when equipped with appropriate structural biases.


% \subsection{GraphGPS \cite{graphgps}}

% GraphGPS proposes a scalable architecture that combines local message passing with global transformer attention. The paper highlights that successful graph transformers typically require three key components: structural positional encodings, local neighborhood aggregation, and global attention.

% In GraphGPS, structural encodings such as Laplacian eigenvectors or random walk positional features are added to node representations to provide information about graph topology. These encodings play a similar role to the spatial and centrality encodings introduced in Graphormer.

% The results show that combining structural encodings with transformer attention leads to strong performance on large-scale graph benchmarks. This work reinforces the central argument of Graphormer that the effectiveness of transformers for graph learning depends on how well graph structure is encoded.


%------------------------------------------------

% \section{Proposed Approach}

% Briefly describe the approach you intend to explore.

% This may include:

% \begin{itemize}
% \item Model or system architecture
% \item Dataset to be used
% \item Evaluation methodology
% \item Expected challenges
% \end{itemize}

%------------------------------------------------

% \section{Conclusion}

% Summarize the problem and the proposed direction of the project.
% Mention the next steps planned for completing the work.

%------------------------------------------------
\bibliographystyle{plain}
\bibliography{references}


\end{document}
